\begin{apendicesenv}

\partapendices

\chapter{Especificação da API REST}

Este apêndice documenta o contrato de interface da API Backend for Frontend (BFF)
desenvolvida. A comunicação é realizada via HTTP utilizando JSON para troca de dados.

\section{Endpoint de Chat e Anamnese}

\textbf{Rota:} \texttt{POST /api/message}

\textbf{Descrição:} recebe mensagem do usuário, gerencia histórico da sessão em memória,
aciona os modelos de IA e retorna resposta conversacional com entidades extraídas.

\subsection{Requisição}

\textbf{Content-Type:} \texttt{application/json}

\textbf{Body de exemplo:}

\begin{verbatim}
{
  "userId": "paciente_12345",
  "text": "Estou sentindo dor de cabeca muito forte",
  "context": "com enjoo desde ontem"
}
\end{verbatim}

\subsection{Resposta}

\textbf{Status Code:} \texttt{200 OK}

\textbf{Body de exemplo:}

\begin{verbatim}
{
  "reply": "Sinto muito. Você poderia descrever se essa dor é pulsante?",
  "entities": [
    {
      "label": "B-Sign_symptom",
      "text": "dor de cabeça",
      "score": 0.98
    },
    {
      "label": "B-Sign_symptom",
      "text": "enjoo",
      "score": 0.95
    }
  ]
}
\end{verbatim}

\subsection{Códigos de Erro}

\begin{itemize}
\item \texttt{400 Bad Request}: quando \texttt{userId} ou \texttt{text} não são fornecidos.
\item \texttt{500 Internal Server Error}: falha na comunicação com a API de inferência.
\end{itemize}

\chapter{Códigos-fonte do Orquestrador (BFF)}

O código do servidor Node.js utilizado no protótipo funcional está no arquivo:
\path{medico_digital/backend/src/server.js}.

Este módulo implementa:

\begin{itemize}
\item Inicialização do servidor Express e middlewares.
\item Gestão de contexto em memória por \texttt{userId}.
\item Chamada ao modelo generativo para condução da anamnese.
\item Chamada ao modelo NER para extração de entidades clínicas.
\item Sanitização da resposta textual e retorno unificado em JSON.
\end{itemize}

\chapter{Artefatos de Teste Exploratórios}

Este apêndice consolida evidências iniciais da avaliação técnica realizada sobre o protótipo
funcional.

\section{Cenários Executados}

\begin{table}[htb]
\centering
\caption{Resumo dos cenários exploratórios executados.}
\label{tab:apendice-testes}
\begin{tabular}{p{2.0cm}p{6.0cm}p{5.0cm}}
\hline
Caso & Objetivo & Resultado observado \\
\hline
CT01 & Validar condução de pergunta única por interação & Comportamento aderente na maior parte das execuções \\
CT02 & Verificar extração de múltiplas entidades clínicas & Identificação satisfatória dos sintomas centrais \\
CT03 & Verificar reconhecimento de medicamento informado pelo paciente & Medicamento extraído corretamente no retorno estruturado \\
CT04 & Observar comportamento em sentenças com negação & Limitação identificada e registrada para tratamento posterior \\
CT05 & Medir latência ponta a ponta do fluxo principal & Faixa observada entre 2.5 s e 4.0 s em ambiente local \\
\hline
\end{tabular}
\end{table}

\section{Síntese das Evidências}

Os testes exploratórios não substituem uma validação com usuários ou um protocolo estatístico
formal. Ainda assim, servem como evidência inicial de coerência arquitetural, aderência ao
fluxo conversacional proposto e viabilidade de integração entre o componente generativo e o
componente extrativo.

\end{apendicesenv}
